\subsection{Uniformly controlled complex Hessians}
\label{sec:gaud-hessian-adjoint-bounds}

The preceding construction bounds the solution form and its Gauduchon multiplier by controlling derivatives of the logarithmic density. A bound on the solution form itself leads to a different estimate. If $\Omega=\omega+\ddc v$, then $\mathrm d\Omega=\mathrm d\omega$, so the exterior derivative of the cofactor $\Omega^{n-1}$ depends only on the fixed torsion and the size of $\Omega$. This permits uniform bounds for the multiplier even when derivatives of the varying complex Hessian are unbounded.

For $n=2$, every Hermitian form has positive lower mass by \cref{cor:lower-mass-dimension}, so the critical integrability in \cref{thm:gaud-hessian-adjoint-endpoint} below, as in \cref{thm:gaud-local-adjoint-endpoint}, also follows from \cref{thm:gaud-determinant-endpoint,cor:gaud-adjoint-lq}; in that dimension the point of these theorems is the pointwise form of their bounds. We first give the real elliptic estimate used below. A cube denoted by $tQ$ has the same center as $Q$ and $t$ times its side length, and $\fint_Q$ denotes its Euclidean average.

\Needspace{12\baselineskip}
\begin{lemma}\label{lem:gaud-divergence-harnack}
Let $Q\subset\R^m$, $m>2$, be a cube. Suppose that $u$ is a smooth strictly positive solution on a neighborhood of $3Q$ of
\[
 \operatorname{div}(A\nabla u+Bu)=0,
 \qquad \lambda I\leq A\leq\Lambda I,\qquad |B|\leq K,
\]
where $A$ is a smooth symmetric matrix field and $B$ is a smooth vector field. Then
\[
 \sup_{Q/2}u\leq C\inf_{Q/2}u,
\]
where $C$ depends only on $m,\lambda,\Lambda,K$ and the side length of $Q$. Derivatives of $A$ and $B$ do not enter $C$.
\end{lemma}
\begin{proof}
We first estimate positive and negative powers, then compare their small moments by an estimate for $\log u$.

For $p\ne0,1$, test the weak equation by $\eta^2u^{p-1}$, where $\eta$ is a smooth cutoff. The principal term is
\[
 (p-1)\int\eta^2u^{p-2}\langle A\nabla u,\nabla u\rangle.
\]
The other terms contain respectively $2\eta u^{p-1}A\nabla u\cdot\nabla\eta$, $(p-1)\eta^2u^{p-1}B\cdot\nabla u$, and $2\eta u^pB\cdot\nabla\eta$. Taking absolute values after moving the principal term to one side and applying Young's inequality absorbs half of its energy. Thus, for either sign of $p$,
\begin{equation}\label{eq:gaud-harnack-power}
 \int\eta^2|\nabla u^{p/2}|^2
 \leq C_1p^2(1+|p-1|^{-2})
        \int (|\nabla\eta|^2+K^2\eta^2)u^p,
\end{equation}
with $C_1$ depending only on the ellipticity bounds. Each test is legitimate since this individual smooth positive solution is bounded above and away from zero on the cutoff support.

Put $\beta=m/(m-2)>1$ and choose $p_0=\beta^{-k-1/2}$, with $k$ to be fixed later. Apply the Euclidean Sobolev inequality to $\eta u^{p/2}$ on nested cubes shrinking from $Q$ to $Q/2$. Choose the cutoff gaps proportional to $2^{-j}$ and the powers $p_j=p_0\beta^j$. The positive powers never equal one: the closest two are $\beta^{-1/2}$ and $\beta^{1/2}$. Consequently the constants in \eqref{eq:gaud-harnack-power}, together with the cutoff terms from $\nabla(\eta u^{p/2})$, grow at most polynomially in $p_j$ and exponentially in $j$. Since $\sum_j(j+1)\beta^{-j}<\infty$, their iteration product converges. Negative powers give the same calculation for $u^{-1}$. We obtain
\begin{equation}\label{eq:gaud-harnack-moments}
 \sup_{Q/2}u\leq C_2\left(\fint_Q u^{p_0}\right)^{1/p_0},
 \qquad
 \inf_{Q/2}u\geq C_2^{-1}
                  \left(\fint_Q u^{-p_0}\right)^{-1/p_0}.
\end{equation}
Here $C_2$ also depends on $p_0$ and the fixed cube size.

Testing instead by $\eta^2/u$ gives
\[
 \int\eta^2|\nabla\log u|^2
       \leq C_3\int(|\nabla\eta|^2+K^2\eta^2).
\]
For every subcube $S\subset Q$ of side length $\ell$, a cutoff supported in $2S\subset3Q$, equal to one on $S$, has gradient bounded by a dimensional constant times $1/\ell$. Poincar\'e's inequality and Cauchy--Schwarz therefore give
\[
 \fint_S|\log u-(\log u)_S|
 \leq C_4(1+K\ell)\leq K_0,
\]
where $K_0\geq1$ is uniform and $(\log u)_S$ denotes the mean. We spell out the exponential consequence of this mean oscillation bound.

Write $f=\log u$ and fix $a>2$. Select the maximal dyadic subcubes $S$ of $Q$ on which $\fint_S|f-f_Q|>aK_0$. Their total measure is at most $|Q|/a$. The parent of each selected cube has average at most $aK_0$, hence $|f_S-f_Q|\leq2^maK_0$. Outside their union, differentiation of dyadic averages gives $|f-f_Q|\leq aK_0$ almost everywhere. Repeat this selection inside each chosen cube, subtracting its own mean. The union at generation $j$ has measure at most $a^{-j}|Q|$, and outside generation $j+1$ the accumulated mean changes and last oscillation give
\[
 |f-f_Q|\leq (j2^ma+a)K_0.
\]
The intersection of all generations is null. Choosing $c>0$ so that $\exp(c2^ma)<a$ and summing over the generations yields
\[
 \fint_Q\exp\!\left(\frac{c|f-f_Q|}{K_0}\right)\leq C_5,
\]
with dimensional constants. Now fix $k$ large enough that $p_0<c/K_0$. Multiplying the estimates for the positive and negative exponentials cancels $f_Q$ and gives
\[
 \left(\fint_Q u^{p_0}\right)
 \left(\fint_Q u^{-p_0}\right)\leq C_5^2.
\]
Combining this with \eqref{eq:gaud-harnack-moments} proves the result. Every constant comes from the stated coefficient bounds and the fixed cube.
\end{proof}

\Needspace{13\baselineskip}
\begin{proposition}\label{prop:gaud-hessian-factor-bounds}\index{Gauduchon factor!uniform complex-Hessian bounds}
Fix $0<\delta\leq\Lambda$. For every smooth real $v$ with
\[
 \delta\omega\leq\Omega_v:=\omega+\ddc v\leq\Lambda\omega,
\]
let $\rho_v$ be its Gauduchon multiplier, normalized by
\[
 \ddc(\rho_v\Omega_v^{n-1})=0,
 \qquad \int_X\rho_v\,\mathrm{d}V=1.
\]
There is $C_H\geq1$, depending only on $X,\omega,\mathrm{d}V,\delta,\Lambda$, such that $C_H^{-1}\leq\rho_v\leq C_H$ on $X$ for every such $v$.
\end{proposition}
\begin{proof}
Existence and uniqueness of $\rho_v$ follow from \cref{thm:gauduchon}. Fix a real reference metric on $X$, and put $Q_v=\Omega_v^{n-1}$ and $J_v=\Omega_v^n/\mathrm{d}V$. Since $\mathrm d\Omega_v=\mathrm d\omega$,
\begin{equation}\label{eq:gaud-hessian-cofactor-divergence}
 \mathrm dQ_v=(n-1)\mathrm d\omega\wedge\Omega_v^{n-2}.
\end{equation}
Both $Q_v$ and $\mathrm dQ_v$ have uniformly bounded coefficients in fixed charts.

Define a real tensor $A_v$ and vector field $b_v$ by
\[
 \langle A_v\mathrm d\phi,\mathrm d\psi\rangle\,\mathrm{d}V
       =n\,\mathrm d\psi\wedge\dc\phi\wedge Q_v,
 \qquad
 (b_v\cdot\mathrm d\phi)\,\mathrm{d}V
       =-n\,\dc\phi\wedge\mathrm dQ_v
\]
for smooth real functions $\phi,\psi$. Only the $(1,1)$ part of $\mathrm d\psi\wedge\dc\phi$ contributes to the first identity; it is
\[
 \frac{\mathrm i}{4\pi}
 (\partial\psi\wedge\bar\partial\phi
       +\partial\phi\wedge\bar\partial\psi).
\]
Thus $A_v$ is symmetric and uniformly positive definite, with a uniform upper bound. The second identity and \eqref{eq:gaud-hessian-cofactor-divergence} give a uniform bound for $b_v$.

For $L_{\Omega_v}\phi=n\ddc\phi\wedge Q_v/\Omega_v^n$, Stokes' theorem gives
\begin{align*}
 \int_X\psi J_vL_{\Omega_v}\phi\,\mathrm{d}V
 &=-n\int_X\mathrm d\psi\wedge\dc\phi\wedge Q_v
          +n\int_X\psi\dc\phi\wedge\mathrm dQ_v\\
 &=-\int_X\langle A_v\mathrm d\phi,\mathrm d\psi\rangle\,\mathrm{d}V
          -\int_X\psi b_v\cdot\mathrm d\phi\,\mathrm{d}V.
\end{align*}
For $\psi=\rho_v$ the left side vanishes for every $\phi$, since $\ddc(\rho_vQ_v)=0$. As $A_v$ is symmetric, the smooth function $\rho_v$ therefore solves
\[
 \operatorname{div}_{\mathrm{d}V}(A_v\mathrm d\rho_v+b_v\rho_v)=0.
\]
In each coordinate chart, multiplying the coefficients by the fixed smooth volume density converts this to Euclidean divergence form while preserving uniform ellipticity and a uniform bound for the vector coefficient. \Cref{lem:gaud-divergence-harnack} applies, with $m=2n\geq4$. Cover $X$ by the half cubes $Q/2$ of finitely many coordinate cubes $Q$ for which $3Q$ lies in a chart. Since $X$ is connected, the overlap graph of this cover is connected, and chaining the local estimates gives $\sup_X\rho_v\leq C_H\inf_X\rho_v$. The normalization and $\mathrm{d}V(X)=1$ imply the claimed two-sided bound.
\end{proof}

For a fixed $0<\varepsilon<1$, consider the full family
\begin{equation}\label{eq:gaud-controlled-hessian-family}
 \mathcal U_\varepsilon=
 \{v\in C^\infty(X,\R):\sup_Xv=0,
                    -\varepsilon\omega\leq\ddc v\leq\varepsilon\omega\}.
\end{equation}
With $J_v=\Omega_v^n/\mathrm{d}V$ as in the preceding proof, set $c_v=\int_X\Omega_v^n$ and $F_v=\log(J_v/c_v)$, so that $\int_X\mathrm e^{F_v}\mathrm{d}V=1$. Here the density $F_v$ is read off from a given solution form instead of being prescribed in advance.

\Needspace{14\baselineskip}
\begin{theorem}\label{thm:gaud-hessian-adjoint-endpoint}\index{Adjoint measure!uniform complex-Hessian bounds}\index{Determinant density!exact maximizers with controlled Hessian}
Let
\[
 j_-=(1-\varepsilon)^n\min_X\frac{\omega^n}{\mathrm{d}V},
 \qquad j_+=(1+\varepsilon)^n\max_X\frac{\omega^n}{\mathrm{d}V},
\]
and take $C_H$ from \cref{prop:gaud-hessian-factor-bounds} with $\delta=1-\varepsilon$ and $\Lambda=1+\varepsilon$. For every $v\in\mathcal U_\varepsilon$, the exact maximizing current $T_v\in\mathscr G_\omega$ for $F_v$ is unique. Its determinant density $k_{T_v}$ and the density $r_v=\mathrm d\nu_{T_v,F_v}/\mathrm{d}V$ of its supporting probability satisfy
\begin{equation}\label{eq:gaud-hessian-adjoint-bounds}
 \frac{j_-}{C_Hj_+}\leq r_v\leq\frac{C_Hj_+}{j_-},
 \qquad
 j_+^{-1/n}\frac{j_-}{C_Hj_+}
 \leq k_{T_v}\leq j_-^{-1/n}\frac{C_Hj_+}{j_-}.
\end{equation}
In particular both densities are uniformly bounded at the critical exponent $n/(n-1)$, without an assumption on the lower mass. Moreover,
\[
 \|F_v-F_0\|_\infty
       \leq n\log\frac{1+\varepsilon}{1-\varepsilon},
\]
where $F_0=F_*=\log(\omega^n/(c_0\,\mathrm{d}V))$ is the density for $v=0$, and this normalized family contains infinitely many distinct nonconstant densities on the fixed manifold $X$.
\end{theorem}
\begin{proof}
Use the normalization of \cref{prop:gaud-hessian-factor-bounds}. The equation $\Omega_v^n=c_v\mathrm e^{F_v}\mathrm{d}V$ gives $B_0(F_v)=\log c_v$, so $\Omega_v$ is the solution form for $F_v$. The identification of the maximizing current and of its densities \eqref{eq:gaud-local-adjoint-densities} in the proof of \cref{thm:gaud-local-adjoint-endpoint} uses only that $\Omega_F$ is the solution form of a smooth datum and that $\chi_F>0$ is a Gauduchon multiplier of $\Omega_F$. Applied with $\Omega_v$ and $\rho_v$ in place of $\Omega_F$ and $\chi_F$, and with
\[
 N_v=\int_X\rho_v\Omega_v^n=\int_XJ_v\rho_v\,\mathrm{d}V
\]
in place of $Z$, it shows that $T_v=(\rho_v/N_v)\Omega_v^{n-1}$ is the unique maximizing current and that
\begin{equation}\label{eq:gaud-hessian-adjoint-densities}
 k_{T_v}=\frac{\rho_v}{N_v}J_v^{(n-1)/n},
 \qquad
 r_v=\frac{J_v\rho_v}{N_v}.
\end{equation}
Here $j_-\leq J_v\leq j_+$ and, since $\int_X\rho_v\,\mathrm{d}V=1$, also $j_-\leq N_v\leq j_+$. The bounds for $\rho_v$ in \cref{prop:gaud-hessian-factor-bounds} give the first pair of inequalities in \eqref{eq:gaud-hessian-adjoint-bounds}; the identity $k_{T_v}=J_v^{-1/n}r_v$ gives the second. The critical norms are bounded by the respective upper constants because $\mathrm{d}V$ is a probability volume.

The inequalities
\[
 (1-\varepsilon)^n\leq\frac{J_v}{J_0}\leq(1+\varepsilon)^n,
 \qquad
 (1-\varepsilon)^n\leq\frac{c_v}{c_0}\leq(1+\varepsilon)^n
\]
give the bound for $F_v-F_0$ by subtracting logarithms. To prove nontriviality, choose a nonconstant smooth real $\phi$. For small $|t|$, the potential $v_t=t\phi-\sup_X(t\phi)$ belongs to $\mathcal U_\varepsilon$. If $F_{v_s}=F_{v_t}$, then $v_s$ and $v_t$ solve the same smooth equation with the same normalization, so $v_s=v_t$ by the uniqueness in \cref{thm:tw-smooth}; hence $(s-t)\phi$ is constant and $s=t$. Thus the densities $F_{v_t}$ are pairwise distinct. Since $\int_X\mathrm e^{F_v}\mathrm{d}V=1$ and $\mathrm{d}V$ is a probability volume, the only constant density in the family is zero, so at most one $F_{v_t}$ is constant.
\end{proof}

Uniform bounds for the solution forms need not give derivative bounds, or even equicontinuity, for their normalized logarithmic densities.

\Needspace{13\baselineskip}
\begin{example}\label{ex:gaud-oscillatory-controlled-hessians}\index{Adjoint measure!oscillatory controlled Hessians}
For every fixed $0<\varepsilon<1$, there are $v_j\in\mathcal V_\varepsilon$ with $\|v_j\|_{C^0}\to0$ such that their normalized densities $F_{v_j}$ are pairwise distinct, bounded in $C^0$, and unbounded in $C^1$. This family is not equicontinuous. Nevertheless, both densities of its exact maximizers satisfy the uniform bounds of \cref{thm:gaud-hessian-adjoint-endpoint}.
\end{example}
\begin{proof}
Choose holomorphic coordinates centered at a point $p$, put $x=\operatorname{Re}z_1$, and take a smooth real cutoff $\zeta$ supported in the chart and equal to one near $p$. We may assume $0\leq\zeta\leq1$. On the chart set
\[
 \beta=dx\wedge\dc x=\ddc(x^2/2)
       =\frac{\mathrm i}{8\pi}dz_1\wedge d\bar z_1.
\]
This is a nonzero semipositive form of rank one. For a constant $a>0$ to be chosen, extend
\[
 w_j=\frac a{j^2}\zeta\sin(jx)
\]
by zero outside the chart, and put $v_j=w_j-\sup_Xw_j$. The product rule gives
\begin{align*}
 \ddc w_j={}&-a\zeta\sin(jx)\beta
       +\frac a j\cos(jx)
          (d\zeta\wedge\dc x+dx\wedge\dc\zeta)\\
       &+\frac a{j^2}\sin(jx)\ddc\zeta.
\end{align*}
All fixed forms in this expression have compact support after multiplication by the indicated cutoff or its derivatives. Thus a fixed $C>0$ bounds the operator norm of $\ddc w_j$ relative to $\omega$ by $aC$ for every $j\geq1$. Choose $a$ so small that $aC\leq\varepsilon/2$. Then $v_j\in\mathcal V_\varepsilon$, and $\|v_j\|_{C^0}\leq2a/j^2$.

Write $s=\operatorname{tr}_{\omega(p)}\beta>0$. Since $x(p)=0$ and $\zeta=1$ near $p$, we have $\Omega_{v_j}(p)=\omega(p)$ and
\[
 \partial_x\Omega_{v_j}(p)=\partial_x\omega(p)-aj\beta.
\]
Differentiating the logarithmic determinant, and observing that its normalizing integral is constant in the spatial variable, gives the exact identity
\[
 \partial_xF_{v_j}(p)
   =\partial_x\log\frac{\omega^n}{\mathrm{d}V}(p)-ajs.
\]
These derivatives are pairwise different and their absolute values tend to infinity. Hence the densities are pairwise distinct and their $C^1$ norms are unbounded. Their uniform $C^0$ bound is part of \cref{thm:gaud-hessian-adjoint-endpoint}.

To check failure of equicontinuity, take points $p_j^\pm$ with $x=\pm\pi/(2j)$ and all other real coordinates zero. For large $j$ they lie where $\zeta=1$, and
$\Omega_{v_j}(p_j^\pm)=\omega(p_j^\pm)\mp a\beta$. The integral normalizing $F_{v_j}$ cancels from the difference, so
\[
 F_{v_j}(p_j^-)-F_{v_j}(p_j^+)
       \longrightarrow\log\frac{1+as}{1-as}>0.
\]
Indeed, $\beta^2=0$ gives $(\omega(p)\pm a\beta)^n=(1\pm as)\omega(p)^n$, and $as<1$ by the chosen form comparison. Both points tend to $p$, which proves the failure of equicontinuity. The exact-adjoint bounds still apply because the entire sequence lies in $\mathcal V_\varepsilon$.
\end{proof}


The same estimate applies to the family with $\delta\omega\leq\Omega_v\leq\Lambda\omega$, with $C_H$ taken for these $\delta,\Lambda$ and with $j_-$ and $j_+$ replaced by $\delta^n\min_X(\omega^n/\mathrm{d}V)$ and $\Lambda^n\max_X(\omega^n/\mathrm{d}V)$. Applying the estimate to a prescribed family of densities requires uniformly comparable solution forms.
